\subsection{复射影空间}

采用第六章 § 3 第 1 号中回忆过的记号，我们给出与实射影空间的定义类似的下述定义：

定义 1 。射影空间 \(\mathbf{P}_{n}(\mathbf{C})\) ，带有可能作为 \(\mathbf{C}_{n+1}^{*}\) 的拓扑对等价关系 \(\Delta_{n}(\mathbf{C})\) 的商的那个拓扑，称为 \(n\) 维复射影空间。

射影空间 \(\mathbf{P}_{1}(\mathbf{C})\) 称为复射影直线，\(\mathbf{P}_{2}(\mathbf{C})\) 称为复射影平面。

关于实射影空间的大部分论证，只需作很小的修改即可推广到复射影空间。

首先我们看到，拓扑空间 \(\mathbf{P}_{n}(\mathbf{C})\) 是豪斯多夫的，这由第六章 § 3 第 1 号命题 1 的论证得到，该论证只需把 R 换成 C 便逐字适用。其次，第六章 § 3 第 1 号命题 2 的证明表明 \(\mathbf{P}_{n}(\mathbf{C})\) 是紧的且连通的，并且同胚于球面 \(S_{2n+1}\) （视为嵌入在空间 \(C_{n+1}^{*}\) 中、后者与 \(R_{2n+2}^{*}\) 等同）对 \(\Delta_{n}(\mathbf{C})\) 在此球面上诱导的等价关系的商；唯一的差别是：现在若 \(n \geqslant 0\) ，这个关系的等价类都同胚于圆周 \(S_{1}\) 。

由于这个原因，第六章 § 3 第 1 号的命题 3 对复射影空间没有类似物。

其次，如同第六章 § 3 第 2 号那样可以证明，每个 \(p\) 维射影线性簇（在空间 \(\mathbf{P}_n(\mathbf{C})\) 中）是闭集，同胚于 \(\mathbf{P}_p(\mathbf{C})\) ，并且其余集在 \(\mathbf{P}_n(\mathbf{C})\) 中稠密（当 \(p < n\) 时）。第六章 § 3 第 2 号命题 5 的证明可以原封不动地转用，只需以 \(\mathbf{C}\) 代换 \(\mathbf{R}\) ，它表明（若 \(n \geqslant 0\) ）\(\mathbf{P}_n(\mathbf{C})\) 中一个射影超平面的余集同胚于 \(\mathbf{C}^n\) ，因而每一点都有一个同胚于 \(\mathbf{C}^n\) 的邻域。这个结果使我们可以把复数空间 \(\mathbf{C}^n\) 嵌入复射影空间 \(\mathbf{P}_n(\mathbf{C})\) ，其作法是把 \(\mathbf{C}^n\) 与一个射影超平面的余集等同起来，这个超平面称为“无穷远超平面”（通常是方程为 \(x_0 = 0\) 的那个超平面）。在特殊情形 \(n = 1\) ，无穷远超平面是一个点，而亚历山德罗夫定理表明 \(\mathbf{P}_1(\mathbf{C})\) 同胚于空间 \(\tilde{\mathbf{C}}\) ——它是对局部紧空间 \(\mathbf{C}\) 添加一个记作 \(\infty\) 的“无穷远点”加以紧化而得到的。于是第六章 § 2 第 4 号的命题 4 表明，复射影直线 \(\mathbf{P}_1(\mathbf{C})\) 同胚于球面 \(S_2\) 。

把与第六章 § 3 第 4 号的结果类似的、关于取值于 C 中的函数的那些结果的叙述留给读者。

把空间 \(R^{n+1}\) 视为嵌入在 \(C^{n+1}\) 中（第 1 号）。设 f 是 \(C_{n+1}^{*}\) 到其商空间 \(\mathbf{P}_{n}(\mathbf{C})\) 上的典范映射。子空间 \(f(\mathbf{R}_{n+1}^{*})\) 由 \(\mathbf{P}_{n}(\mathbf{C})\) 中至少具有一组实齐次坐标的那些点组成；我们来证明 \(f(\mathbf{R}_{n+1}^{*})\) 同胚于实射影空间 \(\mathbf{P}_{n}(\mathbf{R})\) ，这将使我们能把空间 \(\mathbf{P}_{n}(\mathbf{R})\) 视为嵌入在 \(\mathbf{P}_{n}(\mathbf{C})\) 中。现在，\(\Delta_{n}(\mathbf{C})\) 在 \(R_{n+1}^{*}\) 上诱导的关系是 \(\Delta_{n}(\mathbf{R})\) ；典范映射 \(\varphi\) 把

\[
\mathbf {R} _ {n + 1} ^ {*} / \Delta_ {n} (\mathbf {R}) = \mathbf {P} _ {n} (\mathbf {R})
\]

映到 \(f(\mathbf{R}_{n+1}^*)\) 上，它是连续的（第一章，§ 3，第 6 号，命题 10）；由于

\(P_{n}(\mathbf{R})\) 是紧的，故 \(\varphi\) 必是同胚（第一章，§ 9，第 4 号，定理 2，推论 2）。

我们也可以证明 \(\varphi\) 是双连续的而无须用到 \(\mathbf{P}_n(\mathbf{R})\) 的紧性，方法是引用第一章 § 3 第 6 号命题 10 的判别法（见习题 3）。

既然每个 \(p + 1\) 维向量子空间（ \(R^{n+1}\) 的）都生成一个 \(p + 1\) 维的复向量子空间（在 \(C^{n+1}\) 中），我们就看到，\(\mathbf{P}_{n}(\mathbf{R})\) 中每个 p 维射影线性簇 V（V 称为实射影线性簇）生成一个 p 维射影线性簇 \(V'\) （ \(\mathbf{P}_{n}(\mathbf{C})\) 中的； \(V'\) 称为复射影线性簇），使得 V 是 \(V'\) 在 \(\mathbf{P}_{n}(\mathbf{R})\) 上的迹。此外，当我们允许变量取复值时，V 的每个（齐次）方程组也是 \(V'\) 的一个（齐次）方程组。

